\chapter{The Common-Mode Path}
\label{ch:cm}

CM is its own circuit, so work it separately from DM. CM current flows out on
both lines together and returns through chassis. The noise source is the
converter's primary switching node, a voltage $V_n$ behind the
primary-to-chassis parasitic capacitance $C_p$.

\begin{figure}[ht]
	\centering
	\begin{circuitikz}[scale=0.85, transform shape]
		\draw (0,0) to[generic, l=$Z_{s,cm}$] (0,3)
		to[L, l=$L_{cm}$, i>^=$i_{cm}$] (4,3)
		to[short, -*] (5.5,3) -- (8,3)
		to[sV, l=$V_n$] (8,1.5)
		to[C, l=$C_p$] (8,0);
		\draw[dashed] (5.5,3) to[C, l=$2C_y$, -*] (5.5,0);
		\draw (8,0) -- (0,0);
		\draw (4,0) node[ground] {};
		\node[font=\small, right] at (4.2,-0.35) {chassis};
		\node[font=\small, align=center] at (0,3.6) {source};
		\node[font=\small, align=center] at (8,3.6) {converter};
	\end{circuitikz}
	\caption{CM equivalent circuit. Both lines are in parallel. Y caps (dashed)
		are optional and sit on the converter side of the choke.}
	\label{fig:cm-circuit}
\end{figure}

\section{What the choke looks like to CM and DM}
\label{sec:cm-winding}

\begin{derivation}
	For CM, each winding carries $i/2$ in the same sense. Each winding's
	voltage is its self inductance plus the coupling from the other winding:
	\[
		v = j\omega L\frac{i}{2} + j\omega M\frac{i}{2} = j\omega L\,i
		\qquad (M \approx L).
	\]
	The windings are in parallel, so the total CM current $i$ sees $L$: the
	datasheet CM inductance per winding.

	For DM, the windings carry $+i$ and $-i$, so the coupled terms subtract:
	$v = j\omega(L - M)i$. Only the leakage $L - M$ is left, which is $L_{lk}$.
\end{derivation}

A realistic choke model, and a CM source with feedthrough capacitors to chassis (both in parallel for CM):
\[
	Z_{ch} = \left(\frac{1}{sL_{cm}} + \frac{1}{R_p} + sC_w\right)^{-1}, \qquad
	Z_{s,cm} = R + sL_{lead} + \frac{1}{s\,(2C_{ft})}.
\]
$R_p$ models core loss and $C_w$ the winding capacitance.

\section{No Y caps: the choke can make things worse}

\begin{derivation}
	With $C_p$'s impedance dominating the loop, the CM current without the choke
	is $V_n/Z_{Cp}$, and with it $V_n/(Z_{Cp} + Z_{L_{cm}})$. The insertion loss
	(the ratio of the two) is
	\[
		IL = \abs*{\frac{Z_{Cp} + Z_{L_{cm}}}{Z_{Cp}}}
		= \abs*{\frac{\frac{1}{j\omega C_p} + j\omega L_{cm}}{\frac{1}{j\omega C_p}}}
		= \abs{1 - \omega^2 L_{cm} C_p}.
	\]
	$IL < 1$ (a loss of attenuation) whenever $\omega^2L_{cm}C_p < 2$, that is,
	for every frequency below $\sqrt{2}\,f_{res}$.
\end{derivation}

\begin{equation}\label{eq:cm-res}
	f_{res} = \frac{1}{2\pi\sqrt{L_{cm}C_p}}.
\end{equation}

\begin{idea}
	Below $f_{res}$ the choke and $C_p$ form a series resonance that the noise
	source drives directly. The choke only starts helping well above it.
	Without Y caps, a bigger choke moves the problem down in frequency rather
	than removing it.
\end{idea}

\section{With Y caps}

Y caps give the CM current a local return. The noise current, forced by the
high impedance of $C_p$, now splits between the Y caps and the choke-plus-source path:
\[
	Z_Y = \frac{1}{s\,(2C_y)}, \qquad
	IL \approx \frac{\abs{Z_{ch}}}{\abs{Z_Y}}, \qquad
	f_{res,Y} = \frac{1}{2\pi\sqrt{L_{cm}\cdot 2C_y}}.
\]
The two Y caps are in parallel for CM. There is still a resonance, now
between $L_{cm}$ and $2C_y$, so check that it lands outside any band you care about.

Many systems also require the power input to be DC isolated from chassis,
often specified as a minimum resistance such as \qty{1}{\mega\ohm}. Y caps are
open at DC, so they keep that, but they do add capacitance and AC leakage to
chassis. Check for any limit on line-to-chassis capacitance, and that the
caps' voltage rating covers any isolation test.

\begin{examplebox}[label=ex:cm]{the CM path of the candidate}
	The CM parts:
	\begin{center}
		\small
		\begin{tabular}{@{}ll@{}}
			\toprule
			Element                    & Value                                                                                           \\
			\midrule
			CM choke                   & \qty{1}{\milli\henry} per winding, $R_p = \qty{3}{\kilo\ohm}$, $C_w = \qty{8}{\pico\farad}$     \\
			Primary-to-chassis $C_p$   & \qty{500}{\pico\farad} (get the real value from the converter vendor, or measure it)           \\
			CM source                  & \qty{20}{\micro\farad} to chassis (feedthroughs in parallel) + \qty{1}{\micro\henry} + \qty{5}{\ohm} \\
			Optional Y caps            & \qty{4.7}{\nano\farad} from each line to chassis, converter side of the choke                  \\
			\bottomrule
		\end{tabular}
	\end{center}

	\textbf{No Y caps.} From \zcref{eq:cm-res},
	\[
		f_{res} = \frac{1}{2\pi\sqrt{(\qty{1}{\milli\henry})(\qty{500}{\pico\farad})}} = \qty{225}{\kHz}.
	\]
	The choke makes CM noise \emph{worse} at every frequency below
	$\sqrt{2} \times 225 = \qty{318}{\kHz}$. At \qty{400}{\kHz} the ideal choke
	gives $IL = \abs{1 - (400/225)^2} = 6.7$~dB, or 5.5~dB once core loss is
	included. The worst case is $-7.5$~dB at \qty{212}{\kHz}, and the choke only
	reaches 34.5~dB by \qty{10}{\MHz}.

	\textbf{With Y caps.} $2C_y = \qty{9.4}{\nano\farad}$. At \qty{400}{\kHz},
	$\abs{Z_{ch}} = \qty{1985}{\ohm}$ and $\abs{Z_Y} = \qty{42.3}{\ohm}$, so
	\[
		IL \approx 20\log_{10}\frac{1985}{42.3} = 33.4~\text{dB}
	\]
	(33.8~dB exact, 60.3~dB at \qty{10}{\MHz}). The resonance moves down to
	\[
		f_{res,Y} = \frac{1}{2\pi\sqrt{(\qty{1}{\milli\henry})(\qty{9.4}{\nano\farad})}} = \qty{51.9}{\kHz},
	\]
	with a worst case of $-18.4$~dB at \qty{50.3}{\kHz}. Above \qty{150}{\kHz}
	the insertion loss stays above 17.5~dB. Y caps buy about 27~dB at the
	switching frequency, as long as nothing in the system cares about CM noise
	near \qty{52}{\kHz}.
\end{examplebox}
